\section{Reactive thermodynamic inputs}\label{sec:selected-reactive-inputs}
The inputs in the following tables define the baseline parameterization for all seven cases and the Case 3C verification; the perturbation studies vary them around this baseline.
The tables give component, pair, association and reaction coefficients; full precision and sources are provided with the code.

\begin{table}[pos=htbp]
\centering\small\setlength{\tabcolsep}{4pt}
\caption{Component coefficients. Segment count $m$ and charge $z_i$ are dimensionless; $\sigma$ and the Born diameter $d^{\mathrm{Born}}$ are in \unit{\angstrom}, $\epsilon/k_B$ in \unit{K}, and molar mass $M$ in \unit{kg\,mol^{-1}}. Water uses $\sigma_{\HtwoO}(T)$ below. Values marked ``joint fit'' were fitted in the authors' amine-thermodynamics study; printed values are rounded, and full precision is given in the parameter file.}
\label{tab:epcsaft-pure-parameters}
\begin{tabular}{@{}lrrrrrrl@{}}
\toprule
Species & $m$ & $\sigma$ & $\epsilon/k_B$ & $z_i$ & $M$ & $d^{\mathrm{Born}}$ & Source \\
\midrule
\ce{CO2} & 2.0729 & 2.7852 & 169.21 & 0 & 0.04401 & -- & \cite{Gross2001} \\
\ce{MEA} & 3.0353 & 3.0435 & 277.174 & 0 & 0.06108 & -- & \cite{Baygi2015} \\
\ce{H2O} & 1.252754 & $\sigma_{\HtwoO}(T)$ & 432.3788 & 0 & 0.01801528 & -- & Joint fit \\
\ce{MEAH+} & 1 & 5.010868 & 179.5727 & 1 & 0.06209 & 5.010868 & Joint fit \\
\ce{MEACOO-} & 1 & 5.200200 & 203.8710 & -1 & 0.10408 & 5.200200 & Joint fit \\
\ce{HCO3-} & 1 & 2.9296 & 70 & -1 & 0.0610168 & 3 & \cite{Held2014} \\
\ce{CO3^{2-}} & 1 & 2.4422 & 249.26 & -2 & 0.06001 & 3 & \cite{Held2014} \\
\ce{H3O+} & 1 & 3.4654 & 500 & 1 & 0.01902 & 1.218 & \cite{Held2014,Figiel2025} \\
\ce{OH-} & 1 & 2.0177 & 650 & -1 & 0.01701 & 3.0811 & \cite{Held2014,palascakAccurateExperimentalValues2004} \\
\bottomrule
\end{tabular}
\end{table}
\begin{table}[pos=htbp]
\centering\small
\caption{Ionic diameters in \unit{\angstrom}. The packing diameter is $0.88\sigma_i$ and the Debye--H\"uckel diameter is $\sigma_i$. The two amine-ion Born diameters equal their segment diameters. The \ce{HCO3-} and \ce{CO3^{2-}} Born diameters were set in the amine-thermodynamics study, the \ce{H3O+} value is the hydrogen-ion value of \cite{Figiel2025}, and the \ce{OH-} value was calculated from the hydration free energy of \cite{palascakAccurateExperimentalValues2004}.}
\label{tab:reactive-ion-diameters}
\begin{tabular}{@{}lrrr@{}}
\toprule
Ion & $d_{\mathrm{packing}}$ & $d_{\mathrm{DH}}$ & $d_{\mathrm{Born}}$ \\
\midrule
\ce{MEAH+} & 4.409564 & 5.010868 & 5.010868 \\
\ce{MEACOO-} & 4.576176 & 5.200200 & 5.200200 \\
\ce{HCO3-} & 2.578048 & 2.9296 & 3 \\
\ce{CO3^{2-}} & 2.149136 & 2.4422 & 3 \\
\ce{H3O+} & 3.049552 & 3.4654 & 1.218 \\
\ce{OH-} & 1.775576 & 2.0177 & 3.081077 \\
\bottomrule
\end{tabular}
\end{table}
\begin{table}[pos=htbp]
\centering\small\setlength{\tabcolsep}{4pt}
\caption{Nonzero binary interaction coefficients, $k_{ij}(T)=k_{ij}(\qty{313.15}{K})+b_{ij}\,(1/T-1/\qty{313.15}{K})$. The two water--carbonate-ion slopes are instead linear, $k_{ij}(T)=k_{ij}(\qty{313.15}{K})+a_{ij}\,(T-\qty{313.15}{K})$, with $a_{ij}$ in \unit{K^{-1}}. All unlisted unlike-charge and neutral pairs have $k_{ij}=0$, and dispersion between same-sign ions is excluded \cite{Figiel2025}.}
\label{tab:epcsaft-binary-parameters}
\begin{tabular}{@{}lrrl@{}}
\toprule
Pair & $k_{ij}(\qty{313.15}{K})$ & Slope & Source \\
\midrule
\ce{CO2}--\ce{H2O} & -0.0464169 & $b_{ij}=\qty{-35.6651}{K}$ & \ce{CO2}--water pressure fit \\
\ce{MEA}--\ce{H2O} & -0.0883119 & $b_{ij}=\qty{-5.38292}{K}$ & Joint fit \\
\ce{H2O}--\ce{MEAH+} & -0.245955 & $b_{ij}=\qty{108.624}{K}$ & Joint fit \\
\ce{H2O}--\ce{MEACOO-} & -0.166058 & -- & Joint fit \\
\ce{H2O}--\ce{HCO3-} & 0.170693 & $a_{ij}=0.00393119$ & Sodium-salt fit \\
\ce{H2O}--\ce{CO3^{2-}} & -0.970553 & $a_{ij}=0.00196313$ & Sodium-salt fit \\
\ce{MEAH+}--\ce{HCO3-} & 1 & -- & Joint fit, upper bound \\
\ce{MEAH+}--\ce{MEACOO-} & -1 & -- & Joint fit, lower bound \\
\ce{H2O}--\ce{H3O+} & 0.25 & -- & \cite{Held2014} \\
\ce{H2O}--\ce{OH-} & -0.25 & -- & \cite{Held2014} \\
\bottomrule
\end{tabular}
\end{table}

The segment mixing rules are $\sigma_{ij}=(\sigma_i+\sigma_j)/2$ and $\epsilon_{ij}=(1-k_{ij})\sqrt{\epsilon_i\epsilon_j}$.
The \ce{CO2}--water coefficients were fitted to \ce{CO2}--water pressures \cite{kiepeExperimentalDeterminationPrediction2002} before the joint fit, and the two water--carbonate-ion coefficients come from a fit to sodium-salt solutions in the same study.
All solvation factors are one except water, whose factor is 1.5 \cite{Figiel2025}.
The water diameter follows the correlation form of Held et al.\ \cite{Held2008} with a constant fitted in the joint fit, with $T$ numerically in kelvin:
\begin{equation}
 \sigma_{\HtwoO}(T)=2.780433+10.11e^{-0.01775T}-1.417e^{-0.01146T}\quad(\unit{\angstrom}).
\end{equation}
% B2: selected dielectric and Born equations
The bulk relative permittivity follows the density- and temperature-dependent form of Schreckenberg et al.\ \cite{Schreckenberg2014} with the solvent mixing of Yang et al.\ \cite{Yang2024},
\begin{equation}
\varepsilon_{r,\mathrm{bulk}}=1+\rho_m x_{\mathrm{solv}}\,d_{V,\mathrm{mix}}\left(\frac{d_{T,\mathrm{mix}}}{T}-1\right),\qquad
d_{V,\mathrm{mix}}=\frac{\sum_{s}x_s d_{V,s}}{x_{\mathrm{solv}}},\quad
d_{T,\mathrm{mix}}=\frac{\sum_{s}x_s d_{T,s}}{x_{\mathrm{solv}}},\label{eq:solvent-permittivity}
\end{equation}
where $\rho_m$ is the molar density, the sums run over the neutral species \ce{CO2}, MEA and water, and $x_{\mathrm{solv}}=\sum_s x_s$.
Water uses $d_V=\qty{0.3777}{L\,mol^{-1}}$ and $d_T=\qty{1403}{K}$ \cite{Yang2024}.
The MEA values, $d_V=\qty{0.730987}{L\,mol^{-1}}$ and $d_T=\qty{1019.080}{K}$, were fitted to MEA--water permittivities \cite{agieienkoWhatGasStream2024a}; \ce{CO2} is given $d_V=\qty{e-8}{L\,mol^{-1}}$ and $d_T=\qty{1403}{K}$, so it adds no polarization.
The ionic-region relative permittivity is $\varepsilon_{r,\mathrm{ion}}=8$.
The adopted Born expression includes solvation-shell and dielectric-saturation terms \cite{Figiel2025,Figiel2026Correction}.
With species solvation factors $s_i$ as specified above, the mixture factor and outer ion diameters are
\begin{align}
\overline s&=\sum_i x_i^{\mathrm{true}}s_i,&
d_i^{\mathrm{out}}&=d_i^{\mathrm{Born}}
\left[1+\frac{\overline s-1}{|z_i|}\right],\qquad z_i\ne0,\label{eq:born-shell}\\
D_i^{\mathrm{bulk}}&=\frac{1}{d_i^{\mathrm{out}}},&
D_i^{\mathrm{ion}}&=\frac{1}{d_i^{\mathrm{Born}}}-\frac{1}{d_i^{\mathrm{out}}}.\label{eq:born-diameters}
\end{align}
All nine positive solvation factors enter $\overline s$, so here $\overline s=1+0.5x_{\HtwoO}^{\mathrm{true}}$.
In \cref{eq:epcsaft-born}, the two regions use $D_i^{\mathrm{bulk}}$ with $\varepsilon_{r,\mathrm{bulk}}$ and $D_i^{\mathrm{ion}}$ with $\varepsilon_{r,\mathrm{ion}}=8$; diameters are converted to metres in that equation.
Both extension coefficients are one for these inputs, and no polar term is active.

% B4: association topology
Each neutral species has one donor site $a$ and one acceptor site $b$; only the site pairs below are active.
\begin{table}[pos=htbp]
\centering\small
\caption{Active association site pairs. Energy divided by $k_B$ is in \unit{K}; association volume is dimensionless. The MEA values are from \cite{Baygi2015} and the water values from the joint fit. The reciprocal \ce{CO2}--water pairs are induced-association inputs \cite{Schick2023}, with half the water energy and the water volume \cite{Kleiner2007}; the MEA--water cross pairs follow the Wolbach--Sandler combining rule \cite{wolbachUsingMolecularOrbital1998a}.}
\label{tab:reactive-association}
\begin{tabular}{@{}lll@{}}
\toprule
Site pair & $\epsilon^{\mathrm{assoc}}/k_B$ & $\kappa^{\mathrm{assoc}}$ \\
\midrule
MEA:$a$/MEA:$b$ & 2586.3 & 0.0363054 \\
Water:$a$/water:$b$ & 1756.650 & 0.0376197 \\
\COtwo:$a$/water:$b$ & 878.3249 & 0.0376197 \\
\COtwo:$b$/water:$a$ & 878.3249 & 0.0376197 \\
MEA:$a$/water:$b$ & Combining rule & Combining rule \\
MEA:$b$/water:$a$ & Combining rule & Combining rule \\
\bottomrule
\end{tabular}
\end{table}

The MEA--water association energy is the arithmetic mean of the two self-association energies, and its volume is $\sqrt{\kappa_i\kappa_j}$ multiplied by the size factor $\bigl[\sqrt{\sigma_i\sigma_j}/\sigma_{ij}\bigr]^3$.
The MEA self-association volume is the value of Fakouri Baygi and Pahlavanzadeh \cite{Baygi2015}, converted from their $d_{ij}^3$ convention to the $\sigma_{ij}^3$ convention used here.
The parameter set is the MEA set of the authors' amine-thermodynamics study.
That joint fit adjusted the water segment, dispersion and association coefficients, the water-diameter constant, the MEAH$^+$ and MEACOO$^-$ segment diameters and dispersion energies, and the five ion and MEA--water interaction coefficients labelled ``joint fit'' above.
Its objective combined \ce{CO2} partial pressures \cite{Hilliard2008,Jou1995}, species mole fractions \cite{Bottinger2008,matinFacileMethodDetermination2012,Jakobsen2005}, integral absorption heats \cite{HemaBhavya2025}, MEA--water bubble temperatures and vapor compositions \cite{caiBinaryIsobaricVaporLiquid1996}, solution densities \cite{Amundsen2009,Hartono2014} and solution heat capacities \cite{Hilliard2008}.
The absorber data of this work were not used in that fit.


The formulation couples the five reactions in \cref{eq:reactive-r1,eq:reactive-r2,eq:reactive-r3,eq:reactive-r4,eq:reactive-r5} to the liquid EOS, conserved quantities and electroneutrality.
The EOS supplies activities and fugacity coefficients, and the temperature-dependent reaction constants specify chemical equilibrium.
All five constants are fixed literature correlations of the form
\begin{equation}
\ln K_r(T)=A_r+\frac{B_r}{T}+C_r\ln T+D_rT,\label{eq:reaction-constant-form}
\end{equation}
with $T$ numerically in kelvin and the coefficients in \cref{tab:reaction-constant-coefficients}.
\begin{table}[pos=htbp]
\centering\small\setlength{\tabcolsep}{4pt}
\caption{Reaction-constant coefficients of \cref{eq:reaction-constant-form} on the source basis of each correlation. The source ranges are those stated for each correlation.}
\label{tab:reaction-constant-coefficients}
\begin{tabular}{@{}lrrrrll@{}}
\toprule
Reaction & $A_r$ & $B_r$ (\unit{K}) & $C_r$ & $D_r$ (\unit{K^{-1}}) & Basis & Source, range (\unit{K}) \\
\midrule
R1 & 132.899 & -13445.9 & -22.4773 & 0 & Mole fraction & \cite{Austgen1991}, 273.15--498.15 \\
R2 & 235.4815 & -12092.1 & -36.7816 & 0 & Molality & \cite{Austgen1991}, 273.15--498.15 \\
R3 & 216.049 & -12431.7 & -35.4819 & 0 & Mole fraction & \cite{Austgen1991}, 273.15--498.15 \\
R4 & 2.151 & -1545.3 & 0 & 0 & Molality & \cite{Tong2012}, 293.15--323.15 \\
R5 & -0.890870 & -6166.116 & 0 & -0.000984816 & Molality & \cite{batesAcidicDissociationConstant1951}, 273.15--323.15 \\
\bottomrule
\end{tabular}
\end{table}
The R5 coefficients are the Bates and Pinching relation $-\log_{10}K_5=2677.91/T+0.3869+0.0004277T$ written in natural logarithms.
Each source constant refers to infinite dilution in water on its stated basis, with standard molality \qty{1}{mol\,kg^{-1}} for the molality basis.
The thermodynamic calculation converts it to the EOS fugacity basis by evaluating the pure-water solvent reference at the current temperature and the reference pressure \qty{100000}{Pa}, and adds the resulting reference-transfer and pressure offsets to $\ln K_r$.
The effective constants therefore depend on the EOS and reference state as well as on these correlations.
The common calculation range is \qtyrange{293.15}{393.15}{K}; R4 and R5 are extrapolated above \qty{323.15}{K}.
The enhancement concentration divisor is $s_{\COtwo}=1.04542981654115$.

\clearpage
\begingroup
\makeatletter
\setlength{\@fptop}{0pt}
\makeatother
\begin{table*}[pos=t]
\centering\footnotesize
\caption{Convergence and balance diagnostics for the fit-A calculations. Counts identify the calculations included in each group.}
\label{tab:run-diagnostics}
\setlength{\tabcolsep}{4pt}
\begin{tabularx}{\textwidth}{@{}Xrrrrrr@{}}
\toprule
Study & Initial nodes & Final nodes & Tolerance & Iterations & Max. RMS & Max. BC \\
\midrule
Campaign (6) & 21 & 21--22 & 0.5 & 1--2 & $\leq\num{1.55e-1}$ & $\leq\num{2.78e-17}$ \\
Paired Case 3C (2) & 21--41 & 22--42 & 0.05--0.5 & 2 & $\leq\num{1.23e-1}$ & 0 \\
Thermodynamic (8) & 41 & 41--42 & 0.05 & 1--2 & $\leq\num{5.00e-2}$ & 0 \\
Transport (8) & 41 & 41--42 & 0.05 & 1--2 & $\leq\num{4.98e-2}$ & $\leq\num{4.34e-19}$ \\
Operating (7) & 21 & 22 & 0.5 & 2 & $\leq\num{1.81e-1}$ & $\leq\num{1.74e-18}$ \\
Timing (3) & 21 & 22 & 0.5 & 2 & $\leq\num{1.23e-1}$ & 0 \\
Additional Case 3C refinements & \multicolumn{6}{l}{Original run inputs unavailable} \\
\bottomrule
\end{tabularx}
\par\medskip
\begin{tabularx}{\textwidth}{@{}Xrrrr@{}}
\toprule
Study & Component & Charge & \multicolumn{2}{c}{Net-enthalpy-flow range (\unit{W})} \\
 & \multicolumn{2}{c}{Discrepancy (\unit{mol.s^{-1}})} & 101 positions & 10,001 positions \\
\midrule
Campaign (6) & $\leq\num{1.71e-12}$ & $\leq\num{3.49e-12}$ & 128.00--589.60 & -- \\
Paired Case 3C (2) & $\leq\num{1.43e-12}$ & $\leq\num{2.30e-12}$ & 267.57--286.44 & -- \\
Thermodynamic (8) & $\leq\num{1.49e-12}$ & $\leq\num{3.55e-12}$ & 265.94--294.85 & -- \\
Transport (8) & $\leq\num{1.48e-12}$ & $\leq\num{2.30e-12}$ & 267.20--296.04 & -- \\
Operating (7) & $\leq\num{1.59e-12}$ & $\leq\num{2.69e-12}$ & 162.36--378.73 & -- \\
Timing (3) & $\leq\num{1.38e-12}$ & $\leq\num{2.30e-12}$ & 286.44 & -- \\
Additional Case 3C refinements & \multicolumn{4}{l}{Original run inputs unavailable} \\
\bottomrule
\end{tabularx}
\par\smallskip
\begin{minipage}{\textwidth}\footnotesize
RMS is the maximum normalized differential residual. BC is the maximum scaled boundary residual. Discrepancies are rounded upward. Shared baselines occur in more than one group. Thermodynamic rows exclude the reference. Transport includes the factor-one control. Case 1C is excluded until its original equilibrium start can be reproduced. Dense energy diagnostics require the original additional-refinement runs.
\end{minipage}
\end{table*}

\clearpage
\endgroup
